\documentclass[10pt]{article}
\usepackage[T1]{fontenc}
\usepackage[utf8]{inputenc}
\usepackage[french]{babel}
\usepackage{amsmath,amssymb,mathrsfs,mathpazo}
\usepackage{geometry}
\geometry{a4paper,margin=2cm}
\usepackage{graphicx,array,tabularx,tabularray,colortbl,multirow,multicol,booktabs,enumitem,pifont,ulem,fancybox,diagbox,cancel}
\usepackage{pstricks,pstricks-add,pst-plot,pst-tree,pst-node,pst-text,pst-func,pst-math,pst-bspline,pst-3dplot,pst-eucl,pst-solides3d}
\let\pscalculate\undefined
\usepackage{pst-bezier}
\usepackage{tikz,pgfplots,tkz-tab}
\usetikzlibrary{arrows,arrows.meta,patterns,calc,positioning,trees,decorations.markings,shapes.geometric,shapes.misc,shapes.arrows,intersections,decorations.pathmorphing,decorations.pathreplacing,backgrounds,shadows,babel}
\pgfplotsset{compat=1.18}
\usepgfplotslibrary{fillbetween}
\usepackage[np]{numprint}
\usepackage{listings,makecell,pas-tableur,twoopt}
\let\barre\undefined
\usepackage[pstricks]{tablvar}
\definecolor{darkgreen}{RGB}{0,100,0}
\providecommand{\red}{\color{red}}
\providecommand{\blue}{\color{blue}}
\providecommand{\green}{\color{green}}
\providecommand{\black}{\color{black}}
\providecommand{\white}{\color{white}}
\providecommand{\gray}{\color{gray}}
\providecommand{\cyan}{\color{cyan}}
\providecommand{\magenta}{\color{magenta}}
\providecommand{\yellow}{\color{yellow}}
\providecommand{\violet}{\color{violet}}
\providecommand{\orange}{\color{orange}}
\providecommand{\euro}{\texteuro}
\providecommand{\EUR}{\texteuro}
\providecommand{\ieme}{\textsuperscript{e}}
\providecommand{\degres}{\ensuremath{^\circ}}
\providecommand{\degree}{\ensuremath{^\circ}}
\providecommand{\vv}[1]{\overrightarrow{#1}}
\providecommand{\Coord}[1]{#1}
\providecommand{\decosix}{}
\providecommand{\starredbullet}{\ensuremath{\bullet}}
\pagestyle{empty}
\setlength{\parindent}{0pt}
\usepackage{adjustbox}\DeclareUnicodeCharacter{2236}{:}\DeclareUnicodeCharacter{2190}{\ensuremath{\leftarrow}}
\providecommand{\R}{\mathbb{R}}
\providecommand{\N}{\mathbb{N}}
\providecommand{\Z}{\mathbb{Z}}
\providecommand{\Q}{\mathbb{Q}}
\providecommand{\D}{\mathbb{D}}
\providecommand{\C}{\mathbb{C}}
\providecommand{\e}{{\mathrm{e}}}
\providecommand{\vect}[1]{\overrightarrow{#1}}
\providecommand{\vectt}[1]{\overrightarrow{\mathrm{#1}}}
\providecommand{\vvt}[1]{\overrightarrow{\mathrm{#1}}}
\providecommand{\np}[2][]{{#2~#1}}
\providecommand{\vv}[1]{{\overrightarrow{#1}}}
\renewcommand{\d}{{\mathrm{d}}}
\providecommand{\degres}{{^\circ}}
\providecommand{\degre}{{^\circ}}
\providecommand{\blue}{\color{blue}}
\providecommand{\red}{\color{red}}
\providecommand{\enskip}{\quad}
\providecommand{\llbracket}{\lbrack\!\lbrack}
\providecommand{\rrbracket}{\rbrack\!\rbrack}
\providecommand{\euro}{{\text{€}}}
\providecommand{\no}{{\text{n°}}}
\providecommand{\interval}[2]{{[#1;#2]}}
\providecommand{\intervFO}[2]{{[#1;#2[}}
\providecommand{\dotfill}{\ldots}

\providecommand{\ds}{\displaystyle}
\providecommand{\textding}[1]{#1}
\providecommand{\Oij}{\(\mathrm{O};\vec{\imath},\vec{\jmath}\)}
\providecommand{\Oijk}{\(\mathrm{O};\vec{\imath},\vec{\jmath},\vec{k}\)}
\providecommand{\Ouv}{\(\mathrm{O};\vec{u},\vec{v}\)}
\providecommand{\intervFF}[2]{\ensuremath{[#1;#2]}}
\renewcommand{\euro}{\ifmmode\text{\texteuro}\else\texteuro\fi}
\DeclareUnicodeCharacter{00B0}{\ifmmode^{\circ}\else\textdegree\fi}
\DeclareUnicodeCharacter{20AC}{\ifmmode\text{\texteuro}\else\texteuro\fi}
% Énoncés : Éducation nationale. Transcription et corrigés : APMEP, Denis Vergès et contributeurs. Découpage et classement : CoopMaths, Sébastien Lozano, Éric Elter et contributeurs. Adaptation de la présentation : Maths Explorer.
% CC BY-SA : https://coopmaths.fr/www/about/
% https://coopmaths.fr/alea/static/bac/2024/tex/bac_2024_06_sujet1_metropole_devoile_3.tex
\begin{document}
\noindent\textit{APMEP / CoopMaths. CC BY-SA. Présentation adaptée par Maths Explorer.}\par\medskip
\medskip

\begin{center}\textbf{Partie A : étude d'une fonction} \end{center}

On considère la fonction $f$ définie sur $\R$ par 
\[f(x) = x - \ln \left(x^2 + 1\right),\]

où ln désigne la fonction logarithme népérien.

\medskip

\begin{enumerate}
\item On admet que $f$ est dérivable sur $\R$ et on note $f'$ sa fonction dérivée.
	\begin{enumerate}
		\item Montrer que pour tout nombre réel $x$, on a :
\[f'(x) = \dfrac{(x - 1)^2}{ x^2 + 1}.\]

		\item En déduire le sens de variation de la fonction $f$ sur $\R$.
	\end{enumerate}
\item Montrer que pour tout nombre réel $x > 0$, on a :
\[f(x) = x - 2\ln (x) - \ln \left(1 + \dfrac{1}{x^2}\right).\]

\item Calculer la limite de la fonction $f$ en $+\infty$.
\end{enumerate}

\begin{center}\textbf{Partie B : étude d'une suite}
 \end{center}

On considère la suite $\left(u_n\right)$ définie par : 
\[\left\{\begin{array}{r c l}
u_0& =&7\\
u_{n+1} = f\left(u_n\right)& =&u_n - \ln\left(u_n^2 + 1\right)\:\text{pour tout }\:n \in \N
\end{array}\right.\]

\begin{enumerate}
\item Montrer, en utilisant un raisonnement par récurrence, que pour tout entier naturel $n$ :
$u_n \geqslant 0$.
\item Montrer que la suite $\left(u_n\right)$ est décroissante.
\item En déduire la convergence de la suite $\left(u_n\right)$.
\item On note $\ell$ la limite de la suite $\left(u_n\right)$. Déterminer la valeur de $\ell$. \item
	\begin{enumerate}
		\item Recopier et compléter le script ci-dessous écrit en langage Python afin qu'il
renvoie la plus petite valeur de l'entier $n$ à partir de laquelle $u_n \leqslant h$, où $h$ est un nombre réel strictement positif.

\begin{center}
\begin{tabular}{|l|}\hline
from math import log as ln\\
\#permet d'utiliser la fonction ln \\
\#Le Logarithme népérien\\
~\\
def seuil(h):\\
\quad n = 0\\
\quad u =7\\
\quad while \ldots :\\
\qquad n = n+1\\
\qquad u = \ldots\\
\quad return n\\ \hline
\end{tabular}
\end{center}

		\item Déterminer la valeur renvoyée lorsqu'on saisit seuil(0.01) dans la console Python. Justifier la réponse.
		\end{enumerate}
\end{enumerate}

\begin{center}\textbf{Partie C : calcul intégral}\end{center}

\begin{enumerate}
\item Étudier le signe de la fonction $f$ sur $[0~;~+\infty[$.
\item Interpréter graphiquement l'intégrale:

\[I = \displaystyle\int_2^4 f(x)\:\text{d}x.\]

\item On admet dans cette question que, pour tout nombre réel $x \in [2~;~4]$, on a
l'encadrement:

\[0,5x -1 \leqslant f(x) \leqslant 0,25x + 0,25.\]

En déduire l'encadrement :

\[ 1 \leqslant I \leqslant 2.\]
\end{enumerate}

\bigskip

\end{document}
